% Part: first-order-logic
% Chapter: beyond
% Section: introduction

\documentclass[../../../include/open-logic-section]{subfiles}

\begin{document}

\olfileid{fol}{byd}{int}

\olsection{आढावा}

प्रथम-क्रम तर्कशास्त्र हीच काही महत्त्वाची एकमेव तार्किक प्रणाली नाही:
प्रथम-क्रम तर्कशास्त्राचे अनेक विस्तार आणि अनेक भेद आहेत. सामान्यतः एखाद्या
तर्कशास्त्रात भाषेचे आकारिक विनिर्देशन, बहुतेकदा---पण नेहमीच नाही---एक
निगामी प्रणाली आणि बहुतेकदा---पण नेहमीच नाही---एक अभिप्रेत चिन्हार्थमीमांसा
यांचा समावेश असतो. पण या संज्ञेच्या तांत्रिक वापरामुळे एक उघड प्रश्न निर्माण
होतो: सहजप्रज्ञ किंवा तात्त्विक अर्थाने वापरल्या जाणाऱ्या ``तर्कशास्त्र'' या
शब्दाशी प्रथम-क्रम तर्कशास्त्र नसलेल्या तर्कशास्त्रांचा संबंध काय? खाली
वर्णिलेल्या सर्व प्रणाली कोणत्या ना कोणत्या प्रकारच्या युक्तिविचाराचे प्रतिमान
देण्यासाठी रचल्या आहेत; पण त्या तार्किक कशामुळे ठरतात, हे आपण सांगू शकतो का?

याची सोपी उत्तरे मिळत नाहीत. ``तर्कशास्त्र'' हा शब्द वेगवेगळ्या अर्थांनी आणि
वेगवेगळ्या संदर्भांत वापरला जातो; आणि ``सत्य'' या संकल्पनेप्रमाणेच या
संकल्पनेचेही अनेक तात्त्विक भूमिकांतून विश्लेषण झाले आहे. उदाहरणार्थ, कोणती
विधाने अनिवार्यतः सत्य, अनुभवपूर्व सत्य, अतार्किक पदांच्या अर्थनिर्धारणापासून
स्वतंत्रपणे सत्य, त्यांच्या आकारामुळे सत्य किंवा भाषिक संकेतामुळे सत्य आहेत
हे ठरवणे हे तार्किक युक्तिविचाराचे ध्येय आहे, असे कोणी मानू शकेल; आणि यांपैकी
प्रत्येक संकल्पना बरीच स्पष्ट करावी लागते. केवळ गणितात वापरल्या जाणाऱ्या
तर्कशास्त्रापुरते लक्ष मर्यादित केले, तरी त्याची व्याप्ती किती याविषयी फारसे
एकमत नाही. उदाहरणार्थ, फ्रेगे यांनी सुरू केलेल्या {\em तर्कमीमांसावादी}
परंपरेत रसेल आणि व्हाइटहेड यांनी \textit{प्रिन्सिपिया मॅथेमॅटिका} मध्ये
तर्कशास्त्राच्या पायावर गणित विकसित करण्याचा प्रयत्न केला. त्यांची तार्किक
प्रणाली खाली वर्णिलेल्या प्रणालीसारख्या उच्च-प्रकार तर्कशास्त्राचे एक रूप
होती. शेवटी त्यांना बहुतेक निकषांनुसार पूर्णतः तार्किक न वाटणारी स्वयंसिद्धके
समाविष्ट करावी लागली (विशेषतः अनंतत्वाचे स्वयंसिद्धक आणि न्यूनीकरणक्षमतेचे
स्वयंसिद्धक); तरीही उच्च-क्रम युक्तिविचाराची काही रूपे तार्किक म्हणून
स्वीकारली पाहिजेत, असे कोणी मानू शकेल. याउलट, ज्यांच्या सत्ताशास्त्रात
``विधानांना'' संवादाच्या वैध वस्तू म्हणून स्थान नाही असे क्वाइन म्हणतात की,
द्वितीय-क्रम आणि उच्च-क्रम तर्कशास्त्र ही मेंढीचे कातडे पांघरलेल्या संच
उपपत्तीचीच रूपे आहेत; म्हणजेच, विधेयांवर संख्यकीकरण करणाऱ्या प्रणाली पूर्णतः
तार्किक नाहीत.

आत्तासाठी असे तात्त्विक प्रश्न पुढे कधीतरी विचारात घेणेच योग्य ठरेल; आणि
खालील प्रणाली तार्किक असोत वा नसोत, विविध प्रकारच्या युक्तिविचारांची आकारिक
आदर्शीकरणे म्हणून त्यांच्याकडे पाहूया.

\end{document}
